\chapter{Intravitreal Injection}

\section*{Overview}
An intravitreal injection delivers medication directly into the vitreous cavity of the eye to treat retinal diseases, most commonly anti-vascular endothelial growth factor agents for macular degeneration and diabetic eye disease.

\section*{Alternative Names}
Intravitreal injection; anti-VEGF injection; eye injection

\section*{Principle}
A fine needle is passed through the sclera into the vitreous cavity, depositing medication adjacent to the retina. This achieves high intraocular drug concentrations with minimal systemic exposure.

\section*{History}
Intravitreal injections were used for antibiotics and gas in the 1970s. The advent of anti-VEGF therapy in the 2000s made injection the leading treatment for neovascular age-related macular degeneration, diabetic retinopathy, and retinal vein occlusion.

\section*{Why This Test or Procedure Is Ordered}
Performed for neovascular age-related macular degeneration, diabetic macular edema, retinal vein occlusion, uveitis, endophthalmitis, and other retinal diseases, and for corticosteroid delivery.

\section*{Is This a Primary or Secondary Test or Procedure}
Primary first-line therapy for many sight-threatening retinal diseases.

\section*{How This Test or Procedure Is Performed}
Performed as an office procedure. The eye is anesthetized with topical or subconjunctival anesthetic, the surface is sterilized with povidone-iodine, a speculum is placed, and a small-gauge needle is introduced through the pars plana to deliver the drug. The eye is then examined and monitored.

\section*{What Preparation Is Needed}
Usually none beyond stopping topical contact lenses. A discussion of injection frequency and follow-up is essential because treatment is often repeated monthly or every few months.

\section*{Contraindications and Precautions}
Active ocular infection or inflammation is a contraindication. Precaution: aseptic technique is critical because endophthalmitis, while rare, is devastating.

\section*{Risks}
Transient discomfort, subconjunctival hemorrhage, elevated intraocular pressure, floaters, cataract, retinal detachment, and the rare but serious risks of endophthalmitis and retinal tears.

\section*{What Happens After the Test or Procedure}
The patient is monitored briefly, and visual acuity and pressure may be checked. Most people resume normal activity immediately.

\section*{Understanding Your Results}
Vision improvement or stabilization over weeks confirms response. Failure to improve may indicate non-responder status, and the regimen is adjusted.

\section*{Normal Results}
Stabilization or improvement of vision with reduction in retinal fluid on optical coherence tomography.

\section*{Follow-up Tests and Procedures}
Frequent injections and monitoring with visual acuity testing and OCT; long-term surveillance for disease recurrence.

\section*{References}
\begin{enumerate}
  \item MedlinePlus. Intravitreal injection. U.S. National Library of Medicine; 2025.
  \item Current Medical Diagnosis and Treatment. McGraw-Hill; 2025.
\end{enumerate}

\clearpage
